\documentclass[10pt,a4paper]{amsart}
\usepackage[margin=19mm]{geometry}
\usepackage{amsmath,amssymb,mathtools}
\usepackage[scr=rsfs,cal=euler]{mathalfa}
\usepackage{xfrac}
\usepackage[unicode,hidelinks,bookmarks=false]{hyperref}
\newcommand{\ie}{i.e.\ }
\newcommand{\cf}{cf.\ }
\newcommand{\resp}{respectively\ }
\newcommand{\Subsection}{\subsection}
\providecommand{\nobreakdash}{\nobreak\hbox{-}}
\setlength{\emergencystretch}{2em}
\multlinegap0pt
\usepackage{xcolor}
\usepackage{amssymb}
\usepackage[shortlabels]{enumitem}
\newtheorem{theorem}{Theorem}
\title{On the Set-Valued Katugampola Fractional Integral: Properties and Regular Selections}
\author{Parneet Kaur, Rattan Lal,, Ankit Kumar}
\date{Source: arXiv:2609.00702v1 (CC BY 4.0)}
\begin{document}
\maketitle

\noindent\emph{Source and licence.} On the Set-Valued Katugampola Fractional Integral: Properties and Regular Selections. Version arXiv:2609.00702v1 (CC BY 4.0), distributed by arXiv under the Creative Commons Attribution 4.0 International licence, http://creativecommons.org/licenses/by/4.0/, as recorded in the arXiv licence metadata for this exact version and shown on the abstract page https://arxiv.org/abs/2609.00702v1. Full licence notice: \texttt{licenses/arxiv-2609.00702-CC-BY-4.0.txt}.\par\smallskip

\noindent\textit{Editorial scope.} Counted unit: the article's theorem BoundVar (source0.tex lines 325--327). The article's complete original proof of it is delivered with it (source0.tex lines 328--413, one \texttt{proof} environment of 86 lines). No mathematics is written, repaired or re-derived: the statement and the proof are the article's own lines, in the article's own notation, and no retained line is substituted or re-flowed.  Retained uncounted as the source-local context the proof needs: lines 226--228; lines 246--248.  Nothing was omitted from any retained span, so the package carries no omission record.  No retained line contains a citation, so the wrapper carries no bibliography and no bibliography entry.  No retained line contains a figure environment, an image-inclusion command or drawing code, so no image is delivered and the package declares no asset dependency.  Editorial formatting only, in addition: an AMS \texttt{amsart} preamble that restates the article's own declarations and macros used by the retained lines, namely the environments \texttt{theorem} on separate counters, together with the standard packages those lines need.  The retained environments print in this document as Theorem 1 in order of appearance, whereas the article prints the same items with the numbering of its own full document.  The complete native source of the article as distributed in the arXiv e-print for this exact version is bundled unchanged as \texttt{sources/209052/source0.tex}.
 \begin{theorem} \label{convex}
     Let the set-valued map $F:\mathcal{I}\to\mathbb{R}$ be such that $F(x)$ is convex for every $x\in \mathcal{I}$. Then for every $\nu > 0$ and $ \rho \neq -1 $ the Katugampola fractional integral of F preserves convexity, that is, $_a^\rho\mathfrak{D}^\nu F(t)$ is convex for all $t\in \mathcal{I}.$
 \end{theorem}

 \begin{theorem} \label{compact,continuity}
    Let the set valued mapping $F:\mathcal{I} \rightrightarrows \mathbb{R}$ with non-empty compact values, Borel measurable and integrably bounded on $\mathcal{I}.$ Then, for each $t\in \mathcal{I},$ the set valued Katugampola integral is non empty and compact and the mapping  to $_a^\rho\mathfrak{D}^\nu F(t)$ preserves the continuity on $\mathcal{I}$, for $\rho \neq -1, \nu >0,$ in the Hausdorff metric $H_d$. 
\end{theorem}

 \begin{theorem} \label{BoundVar}
     Let the set-valued mapping $F:\mathcal{I} \rightrightarrows \mathbb{R} $ with non empty compact convex values, Borel measurable and integrably bounded on $\mathcal{I}.$ Assume that $F$ is of bounded variation on $\mathcal{I}$ with respect to $H_d$ and that $\rho \neq -1, \nu >1$ Then, Katugampola integral $t\in \mathcal{I},$ the set $_a^\rho\mathfrak{D}^\nu F(t)$ is nonempty, compact and convex, and the mapping  to $_a^\rho\mathfrak{D}^\nu F(t)$ preserves the bounded variation on $\mathcal{I}$ in Hausdorff metric, $H_d.$
 \end{theorem}

 \begin{proof}
     Since $F(t)$ is non-empty, compact, Borel-measurable and integrably bounded, by Theorem \ref{compact,continuity}, for each $t\in \mathcal{I},$ the set $_a^\rho\mathfrak{D}^\nu F(t)$ is non-empty, compact and the mapping $_a^\rho\mathfrak{D}^\nu F(t)$ is continuous on $\mathcal{I}$ with respect to $H_d.$ Furthermore, since $F(t)$ is convex,Theorem \ref{convex} gives us that the mapping $_a^\rho\mathfrak{D}^\nu F(t)$ is convex. Thus, it is sufficient to verify that the mapping $_a^\rho\mathfrak{D}^\nu F(t)$ is of bounded variation. Also, for each $t\in \mathcal{I},$ the set $ F(t)$ nonempty, compact and convex subset of $\mathbb{R}$ and hence it must be a closed interval. 
     \[ \underline{f}(t):= \inf F(t)  \text{ and } \overline{f}(t):=\sup F(t) \]
    The set valued $F(t)$ can be written as 
    \[ F(t)= [\underline{f}(t),\overline{f}(t)], \text{ for all } t\in \mathcal{I} \] 
  Consider a partition $\Gamma=\{0\le a=t_1 <t_2 <\ldots < t_n=b \}$ of interval $\mathcal{I}.$ Now, the Hausdorff distance is measured as the 
  \begin{equation*}
      \begin{aligned}
  H_d(F(t_i),F(t_{i-1}))=& H_d([\underline{f}(t_i),\overline{f}(t_i)],[\underline{f}(t_{i-1}),\overline{f}(t_{i-1})])\\
  =&\max \big\{ |\underline{f}(t_i)-\underline{f}(t_{i-1})|, |\overline{f}(t_{i})-\overline{f}(t_{i-1})|\big\}.
   \end{aligned}
  \end{equation*}
  From this 
   \begin{equation*}
       \begin{aligned}
           \sum_{i=1}^{n}|\overline{f}(t_{i})-\overline{f}(t_{i-1})|\le& \sum_{i=1}^nH_d(F(t_i),F(t_{i-1}))  \text{ and also, } \\
           \sum_{i=1}^{n}|\underline{f}(t_i)-\underline{f}(t_{i-1})|\le& \sum_{i=1}^nH_d(F(t_i),F(t_{i-1})).
       \end{aligned}
   \end{equation*}
  If we take the supremum on all partitions $\Gamma,$ we get
  \begin{equation*}
      \begin{aligned}
          V_a^b(\underline{f})\le V_a^b(F)< \infty  \text{ and }  V_a^b(\overline{f})\le V_a^b(F) < \infty
      \end{aligned}
  \end{equation*}
as $F$ is of bounded variation. From the bounded variation of functions $\underline{f}$ and $ \overline{f}$ on $I,$ these are also bounded on $\mathcal{I}.$ For $t\in I$ define 
\begin{equation*}
    \begin{aligned}
        (_a^\rho\mathfrak{D}^\nu \underline{f})(t)=\frac{(\rho+1)^{1-\nu}}{\Gamma(\nu)}\int_a^t(t^{\rho+1}-x^{\rho+1})^{\nu-1}x^{\rho}\underline{f}(x)dx.
    \end{aligned}
\end{equation*}
Similarly, 
\begin{equation*}
    \begin{aligned}
        (_a^\rho\mathfrak{D}^\nu \overline{f})(t)=\frac{(\rho+1)^{1-\nu}}{\Gamma(\nu)}\int_a^t(t^{\rho+1}-x^{\rho+1})^{\nu-1}x^{\rho}\overline{f}(x)dx.
    \end{aligned}
\end{equation*}
Because $\underline{f}$ and $\overline{f}$ are both bounded and integrable, thus the both operators $_a^\rho\mathfrak{D}^\nu \underline{f}$ and $_a^\rho\mathfrak{D}^\nu \overline{f}$ are well defined. We now assert that each of these operator is Lipschitz and therefore, is of bounded variation on $\mathcal{I}$. Now, as the $\underline{f}$ is bounded, there exist a $M>0$ such that $\underline{f}(t)\le M$ for all $t\in \mathcal{I}. $ Also, as for all $t \in (a,b]$ the operator $_a^\rho\mathfrak{D}^\nu(t)$ is differentiable as the following 
\begin{equation*}
    \begin{aligned}
        (_a^\rho\mathfrak{D}^\nu)' \underline{f}(t)=\frac{(\rho+1)^{2-\nu}(\nu-1)}{\Gamma(\nu)}\int_a^tt^\rho(t^{\rho+1}-x^{\rho+1})^{\nu-2}x^{\rho}\underline{f}(x)dx.
    \end{aligned}
\end{equation*}
Hence, 
\begin{equation*}
    \begin{aligned}
        |(_a^\rho\mathfrak{D}^\nu)' \underline{f}(t)|&=\frac{(\rho+1)^{2-\nu}(\nu-1)}{\Gamma(\nu)}\int_a^tt^\rho(t^{\rho+1}-x^{\rho+1})^{\nu-2}x^{\rho}|\underline{f}(x)|dx \\
        &\le \frac{(\rho+1)^{2-\nu}(\nu-1)}{\Gamma(\nu)}M\int_a^tt^\rho(t^{\rho+1}-x^{\rho+1})^{\nu-2}x^{\rho}dx 
    \end{aligned}
\end{equation*}
Put $t^{\rho+1}-x^{\rho+1}=u$ then $x^\rho dx=-\frac{du}{\rho+1}$ and now the $\int u^{\nu-2}du=\frac{u^{\nu-1}}{\nu-1}.$ So, we get 
\begin{equation*}
    \begin{aligned}
        |(_a^\rho\mathfrak{D}^\nu)' \underline{f}(t)|&\le  \frac{M(\rho+1)^{2-\nu}(\nu-1)}{\Gamma(\nu)} \frac{t^\rho(t^{\rho+1}-a
        ^{\rho+1})^{(\nu-1)}}{(\rho+1)(\nu-1)}\\
        &\le \frac{M(\rho+1)^{1-\nu}}{\Gamma(\nu)} b^\rho (b^{\rho+1}-a
        ^{\rho+1})^{(\nu-1)}
    \end{aligned}
\end{equation*}
Hence, the $(_a^\rho\mathfrak{D}^\nu)' \underline{f}(t)$ is bounded on interval $(a,b].$ Moreover, as $ t \to a ^+$ can be continuously extends to $[a,b],$ so derivative is bounded on entire interval $\mathcal{I}.$ So the $(_a^\rho\mathfrak{D}^\nu)\underline{f}(t)$ is Lipschitz on $\mathcal{I}$ with Lipschitz constant 
\[ \frac{M(\rho+1)^{-\nu}}{\Gamma(\nu)} b^\rho (b^{\rho+1}-a
        ^{\rho+1})^{(\nu-1)} \]
        Similarly, we can show for the $(_a^\rho\mathfrak{D}^\nu) \overline{f}(t)$ is also Lipschitz with same Lipschitz constant. Thus, both are functions of bounded variation of $\mathcal{I}.$\\
        We now established that for every $t\in \mathcal{I}$
    \begin{equation} \label{interval}
        \begin{aligned}  
     _a^\rho\mathfrak{D}^\nu F(t)=[(_a^\rho\mathfrak{D}^\nu) \underline{f}(t),(_a^\rho\mathfrak{D}^\nu)\overline{f}(t) ]. 
       \end{aligned}
    \end{equation}  
        Let $f$ is a integrable selection of the set-valued mappings $F.$ Then, for each $t\in \mathcal{I},$ it satisfies 
        \[ \underline{f}(t)\le f(t) \le \overline{f}(t).\] 
        Since the kernel \[\frac{(\rho+1)^{1-\nu}(t^{\rho+1}-x^{\rho+1})^{\nu-1}x^{\rho}}{\Gamma(\nu)}\] is non negative for all $x\in [a,t],$ integrating the above inequality yields $ (_a^\rho\mathfrak{D}^\nu) \underline{f}(t) \le _a^\rho\mathfrak{D}^\nu f(t)\le  (_a^\rho\mathfrak{D}^\nu) \overline{f}(t).$ 
        This shows that every element of $ _a^\rho\mathfrak{D}^\nu F(t) $ lies within the  interval of (\ref{interval}). Clearly, both $\underline{f}$ and $\overline{f}$ are selections of $F.$ Therefore, their corresponding fractional integral $(_a^\rho\mathfrak{D}^\nu \underline{f})(t)$ and $(_a^\rho\mathfrak{D}^\nu \overline{f})(t)$ belong to $_a^\rho\mathfrak{D}^\nu F(t).$ Hence the interval in \ref{interval} is contained in $_a^\rho\mathfrak{D}^\nu F(t).$ Combining both the inclusions gives the equality of (\ref{interval}).\\
      Now consider a partition $P=\{a=t_0<t_1<\ldots <t_n=b \}$ of the interval $\mathcal{I}.$ Then for each of $i=1,2\ldots n$ we have 
      \begin{equation*}
          \begin{aligned}
       H_d(_a^\rho\mathfrak{D}^\nu F(t_i),_a^\rho\mathfrak{D}^\nu F(t_{i-1}))=& H_d([_a^\rho\mathfrak{D}^\nu \underline{f}(t_i),_a^\rho\mathfrak{D}^\nu \overline{f}(t_i)],[_a^\rho\mathfrak{D}^\nu \underline{f}(t_{i-1}),_a^\rho\mathfrak{D}^\nu \overline{f}(t_{i-1})]) \\
       =& \max\{|_a^\rho\mathfrak{D}^\nu \underline{f}(t_i)-_a^\rho\mathfrak{D}^\nu \underline{f}(t_{i-1})|,|_a^\rho\mathfrak{D}^\nu \overline{f}(t_i)-_a^\rho\mathfrak{D}^\nu \overline{f}(t_{i-1})| \}
        \end{aligned}
      \end{equation*} 
      Hence, 
      \[ \sum_{i=1}^n H_d(_a^\rho\mathfrak{D}^\nu F(t_i),_a^\rho\mathfrak{D}^\nu F(t_{i-1})) \le \sum_{i=1}^n (|_a^\rho\mathfrak{D}^\nu \underline{f}(t_i)-_a^\rho\mathfrak{D}^\nu \underline{f}(t_{i-1})|+|_a^\rho\mathfrak{D}^\nu \overline{f}(t_i)-_a^\rho\mathfrak{D}^\nu \overline{f}(t_{i-1})| ) \] 
      by taking the supremum over all the partitions $P$ of $\mathcal{I}=[a,b].$ We get 
      \[ V_a^b(_a^\rho\mathfrak{D}^\nu F ) \le V_a^b(_a^\rho\mathfrak{D}^\nu \underline{f}) +V_a^b(_a^\rho\mathfrak{D}^\nu \overline{f})< \infty,\] 
      because both $_a^\rho\mathfrak{D}^\nu \underline{f}$ and $_a^\rho\mathfrak{D}^\nu \overline{f}$ are functions of bounded variations. Therefore, the mapping of the $_a^\rho\mathfrak{D}^\nu F$ is of bounded variation on $\mathcal{I}$ with respect to the Hausdorff metric $H_d.$ This completes the argument.  
 \end{proof}

\end{document}
